\begin{figure}[p]\centering\begin{tikzpicture}[gclplate]\node[anchor=west,font=\sffamily\bfseries\Large,gcllabel] at (-6,3.8){DOUBLE NEGATION IS A BOUNDARY};\draw[GCLWarm,line width=1pt](-6,3.45)--(6,3.45);\node[gcllabel,draw](c) at (-3,1.2){classical source};\node[gcllabel,draw](n) at (3,1.2){negative/CPS target};\draw[gclbluearrow](c)--node[above,gcllabel]{translation}(n);\end{tikzpicture}\caption{\textbf{Plate 24. Double negation as a translation boundary.} Classical reasoning can enter a constructive target through a negative translation rather than literal proof identity. Scope: selected translation only.}\label{plate:double-negation-boundary}\end{figure}
